\begin{figure}[H]
    \centering

    % ── RYSUNEK PRZED KONTRAKCJĄ ──────────────────────────────────────────
    \begin{subfigure}[b]{\linewidth}
        \centering
        \begin{tikzpicture}[
            scale=1.2,
            vtx/.style      = {circle, draw, fill=gray!10, minimum size=0.65cm, inner sep=1pt},
            base/.style     = {circle, draw, fill=orange!40, minimum size=0.65cm, inner sep=1pt},
            free/.style     = {circle, draw, fill=green!25, minimum size=0.65cm, inner sep=1pt},
            matched/.style  = {line width=2pt, blue!70!black},
            unmatched/.style= {thick, black},
            annot/.style    = {draw=none, fill=none, font=\small\itshape, inner sep=1pt}
        ]
            % ── węzły ──
            \node[free] (r) at (0,   0)    {$r$};
            \node[vtx]  (v) at (1.6, 0)    {$v$};
            \node[base] (b) at (3.2, 0)    {$b$};
            \node[vtx]  (x) at (4.2, 1.3)  {$x$};
            \node[vtx]  (y) at (5.8, 0.75) {$y$};
            \node[vtx]  (z) at (5.8,-0.75) {$z$};
            \node[vtx]  (q) at (4.2,-1.3)  {$q$};

            % ── krawędzie łodygi ──
            \draw[unmatched] (r) -- (v);
            \draw[matched]   (v) -- (b);

            % ── krawędzie cyklu (kwiatu) ──
            \draw[unmatched] (b) -- (x);
            \draw[matched]   (x) -- (y);
            \draw[unmatched] (y) -- (z);
            \draw[matched]   (z) -- (q);
            \draw[unmatched] (q) -- (b);

            % ── adnotacja: wierzchołek wolny ──
            \node[annot, above=4pt of r] {wolny};

            % ── adnotacja: wierzchołek bazowy ──
            \node[annot, above left=4pt and -1em of b] {bazowy};

            % ── klamra + etykieta „łodyga" ──
            \draw[decorate, decoration={brace, amplitude=6pt, mirror},
                  thick, gray!70]
                ([yshift=-0.5em]r.south west) -- ([yshift=-0.5em]b.south east)
                node[midway, below=8pt, draw=none, fill=none,
                     font=\small, text=gray!80] {łodyga};

            % ── klamra + etykieta „kwiat" ──
            \draw[decorate, decoration={brace, amplitude=6pt}, thick, gray!70]
                ($(b.west |- x.north) + (0, 0.3)$) -- ($(y.east |- x.north) + (0, 0.3)$)
                node[midway, above=8pt, draw=none, fill=none, font=\small, text=gray!80] {kwiat $B$};

        \end{tikzpicture}
        \caption{Graf przed kontrakcją. Kwiat $B = \{b, x, y, z, q\}$ to cykl długości~5
                 z $k{=}2$ krawędziami skojarzenia (niebieskie).
                 Wierzchołek bazowy $b$ (pomarańczowy) ma obie swoje krawędzie w cyklu
                 poza skojarzeniem.
                 Łodyga $r$--$v$--$b$ łączy kwiat z wolnym wierzchołkiem $r$ (zielony);
                 jej krawędź $vb$ należy do skojarzenia.}
    \end{subfigure}

    \vspace{0.8cm}

    % ── RYSUNEK PO KONTRAKCJI ─────────────────────────────────────────────
    \begin{subfigure}[b]{\linewidth}
        \centering
        \begin{tikzpicture}[
            scale=1.2,
            vtx/.style      = {circle, draw, fill=gray!10, minimum size=0.65cm, inner sep=1pt},
            free/.style     = {circle, draw, fill=green!25, minimum size=0.65cm, inner sep=1pt},
            super/.style    = {circle, draw=orange!70!black, fill=orange!20,
                               minimum size=1.4cm, line width=1.5pt, inner sep=1pt},
            matched/.style  = {line width=2pt, blue!70!black},
            unmatched/.style= {thick, black},
            annot/.style    = {draw=none, fill=none, font=\small\itshape, inner sep=1pt}
        ]
            \node[free]  (r) at (0,   0) {$r$};
            \node[vtx]   (v) at (1.6, 0) {$v$};
            \node[super] (B) at (3.4, 0) {$B$};

            \draw[unmatched] (r) -- (v);
            \draw[matched]   (v) -- (B);

            \node[annot, above=4pt of r] {wolny};
            \node[annot, above=5pt of B] {super-wierzchołek};

        \end{tikzpicture}
        \caption{Graf po kontrakcji kwiatu $B$ do pojedynczego super-wierzchołka.
                 Łodyga pozostaje niezmieniona.}
    \end{subfigure}

    % ── LEGENDA ──────────────────────────────────────────────────────────
    \vspace{0.4cm}
    \begin{tikzpicture}
        \draw[line width=2pt, blue!70!black] (0, 0.4) -- (0.8, 0.4)
            node[right, draw=none, fill=none, font=\small] {krawędź należąca do skojarzenia $M$};
        \draw[thick, black] (0, 0) -- (0.8, 0)
            node[right, draw=none, fill=none, font=\small] {krawędź nienależąca do $M$};
    \end{tikzpicture}

    \caption{Kontrakcja kwiatu w algorytmie Edmondsa.}
    \label{rys:kwiat_kontrakcja}
\end{figure}